\begin{frame}{1.1 Inisialisasi State Target}
    Qubit target $| \psi_{target} \rangle$ (pada register $q_3$) disiapkan dalam superposisi non-trivial untuk menyandikan \textit{eigenvector} awal berdasarkan data historis:
    \begin{equation}
        | \psi_{target} \rangle = \begin{pmatrix} c_0 \\ c_1 \end{pmatrix} = \begin{pmatrix} \cos(0.4996) e^{-i0.1144} \\ \sin(0.4996) e^{i(0.3252 - 0.1144)} \end{pmatrix}
    \end{equation}
    \begin{equation}
        | \psi_{target} \rangle \approx \begin{pmatrix} 0.872 - 0.100i \\ 0.468 + 0.100i \end{pmatrix}
    \end{equation}
    Inisialisasi ini memastikan bahwa sirkuit memulai estimasi dari profil volatilitas yang mendekati kondisi riil pasar.
\end{frame}

\begin{frame}{1.2 Transformasi Walsh-Hadamard dan $|\Phi_0\rangle$}
    Register fase ($q_0, q_1, q_2$) diinisialisasi ke state superposisi \textit{uniform} melalui operator $H^{\otimes 3}$:
    \begin{equation}
        H|0\rangle = \frac{1}{\sqrt{2}}(|0\rangle + |1\rangle) \implies H^{\otimes 3}|000\rangle = \frac{1}{\sqrt{8}} \sum_{k=0}^{7} |k\rangle
    \end{equation}
    State total sistem $|\Phi_0\rangle$ dibentuk melalui \textit{tensor product}:
    \begin{equation}
        |\Phi_0\rangle = (H^{\otimes 3} |000\rangle) \otimes |\psi_{target}\rangle \approx \sum_{k=0}^{7} \left( \frac{c_0}{\sqrt{8}}|k,0\rangle + \frac{c_1}{\sqrt{8}}|k,1\rangle \right)
    \end{equation}
    Hal ini memberikan setiap basis komputasi probabilitas yang sama untuk berinteraksi dengan komponen \textit{eigenvector} melalui gerbang \textit{Controlled-Unitary}.
\end{frame}

\begin{frame}{1.3 Interaksi Controlled-Unitary (CU)}
    Operator $U^{2^j}$ diterapkan menggunakan gerbang \texttt{cu(theta, phi, lambda)} dengan matriks uniter $U_j$:
    \begin{equation}
        U_j = \begin{pmatrix} \cos(\theta_j/2) & -e^{i\lambda_j}\sin(\theta_j/2) \\ e^{i\phi_j}\sin(\theta_j/2) & e^{i(\phi_j+\lambda_j)}\cos(\theta_j/2) \end{pmatrix}
    \end{equation}
    \textbf{Parameter Numerik dari Kode:}
    \begin{itemize}
        \item $U^{2^0}$ ($q_2$): $\theta_0=1.59899, \phi_0=-1.11512, \lambda_0=2.02647$
        \item $U^{2^1}$ ($q_1$): $\theta_1=2.22862, \phi_1=0.513123, \lambda_1=3.65472$
        \item $U^{2^2}$ ($q_0$): $\theta_2=0.797922, \phi_2=-4.53103, \lambda_2=-1.38944$
    \end{itemize}
    Proses ini menghasilkan \textit{phase kickback} yang mengonversi fase \textit{eigenvalue} menjadi amplitudo register fase.
\end{frame}

\begin{frame}{1.4 Perhitungan IQFT dan $|\Phi_{final}\rangle$}
    Untuk ekstraksi nilai fase $\phi$ ke domain basis komputasi $|y\rangle$, diterapkan operator $\mathcal{QFT}^\dagger$:
    \begin{equation}
        |\Phi_{final}\rangle = \left( \mathcal{QFT}^\dagger \otimes I \right) |\Phi_1\rangle = \frac{1}{8} \sum_{y=0}^{7} \sum_{k=0}^{7} e^{2\pi i k (\phi - y/8)} |y\rangle \otimes |\psi_{target}\rangle
    \end{equation}
    \textbf{Interpretasi Hasil Melalui Interferensi:}
    \begin{itemize}
        \item \textbf{Interferensi Konstruktif:} Terjadi jika $\phi = y/8$, menghasilkan amplitudo maksimal (puncak probabilitas) pada state $|y\rangle$.
        \item \textbf{Interferensi Destruktif:} Untuk nilai $y$ lainnya, amplitudo saling meniadakan, menekan probabilitas state non-target.
    \end{itemize}
    Pengukuran akhir pada register fase memberikan estimasi biner fase $\phi \approx y/2^3$.
\end{frame}
